% !TEX root = ../thesis.tex
\chapter{Analytická časť}\label{ch:analytics}

\section{Úvod do analytickej časti}
Analytická časť predstavuje podrobnú technickú analýzu problému a možností jeho riešenia. Cieľom tejto kapitoly je zmapovať súčasnú \glslink{monolith}{monolitickú implementáciu} aplikácie hry Dáma, identifikovať jej silné a slabé stránky, definovať funkčné a nefunkčné požiadavky a navrhnúť architektonické a technologické riešenia pre \glslink{microservice}{mikroservisnú verziu}. Na základe tejto analýzy budú stanovené metriky a metodika porovnania oboch prístupov v implementačnej a vyhodnocovacej časti práce.

Táto časť sa sústreďuje na analýzu moderných technológií, ktoré sú vhodné pre návrh a implementáciu mikroservisnej architektúry. \glslink{springboot}{Spring Boot} umožňuje rýchly vývoj \glslink{restapi}{REST API} a poskytuje robustný ekosystém knižníc pre prácu s databázami, bezpečnosťou a webovou komunikáciou. \glslink{postgresql}{PostgreSQL} poskytuje spoľahlivé ACID transakcie a dlhodobé ukladanie údajov. \glslink{docker}{Docker} a \glslink{dockercompose}{Docker Compose} zjednodušujú izoláciu služieb a lokálne nasadzovanie. \glslink{springcloudgateway}{Spring Cloud Gateway} umožňuje centralizované smerovanie požiadaviek, autentifikáciu a aplikáciu bezpečnostných politík.

Frontend v oboch verziách systému (monolit aj mikroservisy) je implementovaný pomocou \glslink{react}{React} SPA, pričom všetka komunikácia s backendom prebieha cez \glslink{restapi}{REST API} bez server-side renderovania.

Okrem toho je možné teoreticky zvážiť použitie \glslink{kubernetes}{Kubernetes} pre orchestráciu kontajnerov, \glslink{apachekafka}{Kafka}/\glslink{rabbitmq}{RabbitMQ} pre asynchrónnu komunikáciu, \glslink{prometheus}{Prometheus} a \glslink{grafana}{Grafana} pre monitoring a \glslink{elkstack}{ELK Stack} alebo \glslink{jaeger}{Jaeger}/\glslink{zipkin}{Zipkin} pre centralizované logovanie a trasovanie požiadaviek. Táto analýza demonštruje možnosti výberu technológií a ich využitie pri návrhu robustného, škálovateľného a modulárneho systému.

\section{Analýza existujúcej monolitickej implementácie}
\subsection{Popis existujúceho riešenia}
Projekt pozostáva z \glslink{monolith}{monolitickej aplikácie} a novej \glslink{microservice}{mikroservisnej verzie}, obe s frontendom spracúvajúcim herné rozhranie a komunikáciu s backendom. Backend spracováva pravidlá hry, validáciu ťahov a perzistenciu v \glslink{postgresql}{PostgreSQL} tabuľkách \texttt{games}, \texttt{score}, \texttt{comment}, \texttt{rating} cez JPA.

Všetky údaje (nicky, avatary, priebeh hry, výsledky, komentáre a hodnotenia) sa načítavajú cez HTTP požiadavky na backend, ktorý spracováva pravidlá hry, validáciu ťahov a perzistenciu v \glslink{postgresql}{PostgreSQL} tabuľkách \texttt{games}, \texttt{score}, \texttt{comment} a \texttt{rating} pomocou JPA.

Interakcia prebieha medzi dvoma hráčmi na jednom zariadení, pričom front-end \glslink{react}{React} aplikácia zabezpečuje celú vizuálnu časť.

\subsection{Prevádzkové správanie a vlastnosti}
\begin{itemize}
    \item Práca s používateľskými údajmi: frontend načítava všetky potrebné údaje (nicky, avatary, výsledky) cez \glslink{restapi}{REST API}. 
    \item Perzistencia výsledkov: výsledky sú ukladané cez služby monolitu do \glslink{postgresql}{PostgreSQL}.
    \item Jednovláknová konzistencia: všetky operácie bežia v jednom backendovom procese.
    \item Dev/Ops zjednodušenie: monolit je ľahko spustiteľný lokálne — samostatný backend + samostatný frontend.
\end{itemize}

\subsection{Identifikované limity a problémy}
\begin{itemize}
    \item Škálovanie: pri náraste zaťaženia (viac súbežných relácií) je potrebné škálovať celý \glslink{monolith}{monolit}, čo je neefektívne pri potrebe zvýšiť kapacitu len časti funkcionality.
    \item Odolnosť voči chybám: chyba v jednom module môže spôsobiť zlyhanie celej aplikácie.
    \item Vývojová súbežnosť: viacerí vývojári pracujúci na rôznych častiach môžu naraziť na konflikty v kóde a nasadzovaní.
    \item Testovateľnosť: ťažšie izolovať moduly pre samostatné testovanie; integrácia je náročnejšia.
    \item Nasadzovanie: každá zmena vyžaduje redeploy celého systému.
\end{itemize}

\section{Požiadavky na riešenie}
\subsection{Funkčné požiadavky}
\begin{enumerate}
    \item Základná hra Dáma – implementácia pravidiel hry (pohyb, povinné bije, kráľovanie, počítanie skóre).
    \item Správa hráčov – zadanie niku a avatara pri začiatku hry, perzistentné ukladanie výsledkov.
    \item Ukladanie výsledkov – persistencia výsledkov, komentárov a hodnotení v databáze.
    \item Front-end rozhranie – \glslink{react}{React} SPA aplikácia pre spustenie hry, zobrazenie dosky, zápis ťahov a leaderboard.
    \item Lokálny režim – aplikácia musí podporovať lokálnu hru dvoch hráčov na jednom zariadení.
\end{enumerate}

\subsection{Nefunkčné požiadavky}
\begin{enumerate}
    \item Výkon – rýchla odozva pri interakcii s hernou doskou, minimálne latencie pri vykresľovaní ťahov.
    \item Udržiavateľnosť – kód má byť modulárny, dobre testovateľný a dokumentovaný.
    \item Nasadzovanie – možnosť ľahkého spustenia lokálne cez \glslink{dockercompose}{Docker Compose}.
    \item Monitorovanie a logovanie – základné metriky a logy pre vyhodnotenie správania systému.
    \item Bezpečnosť – základné opatrenia na strane API (CORS, input validation).
    \item Rozšíriteľnosť – architektúra musí umožniť jednoduché doplnenie ďalších komponentov (napr. chat, multiplayer, štatistiky).
    \item Sledovanie výkonu – podpora pre integráciu s nástrojmi ako \glslink{prometheus}{Prometheus} a \glslink{grafana}{Grafana}.
\end{enumerate}

\section{Návrh mikroservisnej architektúry}
\subsection{Rozdelenie na služby}
Navrhované rozdelenie systému na samostatné služby (každá ako samostatný \glslink{springboot}{Spring Boot} projekt):
\begin{itemize}
    \item frontend – prezentácia a UI; komunikuje s \glslink{apigateway}{API Gateway}
    \item gateway-service – centralizované smerovanie požiadaviek
    \item user-service – správa používateľov (prezývky, avatary), údaje uložené v \glslink{postgresql}{PostgreSQL}
    \item game-service – logika hry, pravidlá, vykonávanie ťahov
    \item score-service – správa výsledkov
    \item comment-service – správa komentárov
    \item rating-service – správa hodnotení
    \item (voliteľne) analytics-service – agregácia štatistík
\end{itemize}

Každý servis má vlastnú databázovú schému (Database per Service) – čím sa dosahuje autonómia dátovej zodpovednosti.

\subsection{Komunikácia medzi službami}
\begin{itemize}
    \item Synchrónny \glslink{restapi}{REST API}: používa sa pre okamžité volania (napr. frontend → gateway → game-service).
    \item Asynchrónne udalosti (teoreticky): pomocou message brokeru (\glslink{rabbitmq}{RabbitMQ} alebo \glslink{apachekafka}{Apache Kafka}) pre decoupling a eventual consistency.
    \item API kontrakty: každá služba definuje vlastné DTO, čím sa zabezpečuje nezávislosť.
\end{itemize}

\subsection{Správa stavu hry}
Herné pole a logika hry sú spracovávané v pamäti aplikácie pomocou objektov. Po ukončení hry game-service zapíše výsledok do \glslink{postgresql}{PostgreSQL} (\texttt{score}, \texttt{rating}, \texttt{comment}).

\subsection{Konzistencia a transakcie}
Pri oddelení perzistentných služieb sa uprednostňuje eventual consistency. Napríklad, po ukončení hry game-service vyšle udalosť na score-service. Pre zložitejšie viackrokové transakcie sa uvažuje o SAGA pattern (orchestrated alebo choreographed). \glslink{postgresql}{PostgreSQL} zabezpečí ACID tam, kde je nutné.

\subsection{Bezpečnosť a politika prístupu}
\begin{itemize}
    \item Gateway zodpovedá za základné bezpečnostné politiky (CORS, rate-limiting).
    \item Input validation je aplikovaná na každej vrstve API.
    \item Prístup do databáz je nastavený s princípom least privilege.
\end{itemize}

\section{Technologická analýza}
\subsection{Hlavný technologický stack a jeho využitie}
\begin{itemize}
    \item \glslink{springboot}{Spring Boot} – framework pre backend: umožňuje rýchly vývoj \glslink{restapi}{REST API} a spracovanie logiky mikroservisov. V projekte sa používa vo všetkých službách (user-service, game-service, score-service a ďalšie). Poskytuje modulárnu štruktúru, integráciu so Spring Data pre prácu s databázou, Spring Security pre autentifikáciu a autorizáciu a Spring Web pre tvorbu REST endpointov.
    \item \glslink{postgresql}{PostgreSQL} – relačná databáza: slúži na perzistentné uchovávanie dát o hráčoch, skóre, komentároch a hodnoteniach. Používa sa pre user-service a score-service.
    \item \glslink{docker}{Docker} a \glslink{dockercompose}{Docker Compose} – kontajnerizácia: izolujú jednotlivé služby, uľahčujú lokálne spustenie a testovanie celého systému.
    \item \glslink{springcloudgateway}{Spring Cloud Gateway} – \glslink{apigateway}{API Gateway}: centralizované smerovanie požiadaviek, autentifikácia, CORS, rate-limiting.
    \item \glslink{react}{React} – single-page aplikácia. Zabezpečuje interaktívne používateľské rozhranie, renderovanie hernej dosky a komunikáciu cez \glslink{restapi}{REST API}.
\end{itemize}

\subsection{Teoretické rozšírenie a ďalšie možnosti}
\begin{itemize}
    \item \glslink{kubernetes}{Kubernetes} – orchestrácia kontajnerov: umožňuje automatické nasadzovanie, škálovanie a správu životného cyklu kontajnerov v produkčnom prostredí.
    \item \glslink{apachekafka}{Apache Kafka} / \glslink{rabbitmq}{RabbitMQ} – messaging systémy pre asynchrónnu komunikáciu medzi mikroservisami.
    \item \glslink{prometheus}{Prometheus} + \glslink{grafana}{Grafana} – nástroje pre monitoring a vizualizáciu výkonu.
    \item \glslink{elkstack}{ELK Stack} – centralizované logovanie a analýza logov mikroservisov.
    \item \glslink{jaeger}{Jaeger} / \glslink{zipkin}{Zipkin} – nástroje pre distribuované trasovanie požiadaviek.
\end{itemize}

\section{Testovanie a meranie výkonnosti}
\subsection{Testovacia stratégia}
\begin{itemize}
    \item Jednotkové testy: JUnit + Mockito pre každú službu.
    \item Integračné testy: Spring Test + Testcontainers pre \glslink{postgresql}{PostgreSQL}.
    \item Kontraktové testy: overenie API kontraktov medzi službami (napr. pomocou Pact).
    \item End-to-end testy: simulácia používateľského scenára (Selenium alebo \glslink{restapi}{REST} volania).
\end{itemize}

\subsection{Metodika merania výkonu}
Merania budú prebiehať v kontrolovanom prostredí (rovnaký hardware, čisté inštancie). Scenáre: otvorenie hry, vykonanie všetkých možných ťahov, ukončenie hry, zápis výsledkov do databázy. Metriky: čas odozvy pri vykresľovaní ťahov a aktualizácii herného poľa, čas zápisu výsledkov, spotreba CPU a pamäte počas hry. Nástroje: systémové metriky (\glslink{docker}{Docker} stats, htop). Okrem kvantitatívnych ukazovateľov sa zhodnotí aj zložitosť údržby, čitateľnosť kódu a modularita.

\subsection{Očakávania}
\begin{itemize}
    \item Monolit: nižšia latencia pri lokálnej hre, ale ťažšie rozšírenie a modularita.
    \item Mikroservisy: lepšia modularita, možnosť rozširovať samostatné služby, mierny režijný overhead pri lokálnom behu.
\end{itemize}

\section{Porovnanie architektúr – analytické zhrnutie}
\begin{tabular}{lcc}
Kritérium & Monolit & Mikroservisy \\
\hline
Škálovateľnosť & Celý systém & Samostatné služby \\
Výkon & Nižšia latencia & Vyššia modularita \\
Údržba & Ťažkopádna & Jednoduchšie rozdelenie \\
Nasadenie & Jednoduché & Zložitejšie, flexibilnejšie \\
Monitorovanie & Minimálne & Podpora \glslink{prometheus}{Prometheus}/\glslink{grafana}{Grafana} \\
Konzistencia & ACID & Eventual consistency \\
Komunikácia & Interná & \glslink{restapi}{REST} / Messaging \\
Bezpečnosť & Lokálna & Centrálna v Gateway \\
Odolnosť voči chybám & Nízka & Izolované zlyhania
\end{tabular}

\section{Riziká a mitigácia}
\subsection{Hlavné riziká}
\begin{itemize}
    \item Komplexita infraštruktúry – aj pri lokálnom spustení môže byť náročné správne nastaviť viacero mikroservisov cez \glslink{dockercompose}{Docker Compose}.
    \item Konzistencia dát medzi službami – aj keď všetko beží lokálne, je potrebné zabezpečiť správny zápis a načítanie výsledkov, bodov a hodnotení.
    \item Chyby pri implementácii logiky hry – napr. nesprávne vykreslenie herného poľa alebo nesprávne výpočty skóre a ratingov.
    \item Časové obmedzenia implementácie – dokončenie všetkých služieb a testovanie v rámci stanoveného termínu.
\end{itemize}

\subsection{Návrhy mitigácie}
\begin{itemize}
    \item Iteratívna implementácia – najprv MVP (game-service + gateway + frontend) pre jednoduchšie testovanie a rozširovanie.
    \item Základné monitorovanie a logovanie – kontrola správneho zápisu výsledkov, bodov a ratingov.
    \item Testovanie logiky hry a konzistencie dát – overenie správneho vykreslenia herného poľa, skóre a hodnotení.
\end{itemize}

\section{Zhrnutie analytickej časti a premostenie do návrhovej časti}
Analytická časť identifikovala hlavné obmedzenia \glslink{monolith}{monolitickej architektúry} (ťažké škálovanie, slabá modularita, závislosť na session) a navrhla moderné \glslink{microservice}{mikroservisné riešenie} založené na \glslink{springboot}{Spring Boot}, \glslink{postgresql}{PostgreSQL}, \glslink{docker}{Docker}, \glslink{springcloudgateway}{Spring Cloud Gateway} a \glslink{react}{React}.

Okrem praktických komponentov bola analyzovaná aj možnosť využitia technológií ako \glslink{kubernetes}{Kubernetes}, \glslink{apachekafka}{Apache Kafka}, \glslink{rabbitmq}{RabbitMQ}, \glslink{prometheus}{Prometheus}, \glslink{grafana}{Grafana} a \glslink{elkstack}{ELK Stack}, ktoré by umožnili škálovanie a monitorovanie systému v produkčnom prostredí.

Navrhnutá architektúra je modulárna, rozšíriteľná a prispôsobená pre objektívne porovnanie s monolitom podľa výkonových a údržbových metrík. Výsledky tejto analýzy budú tvoriť základ pre implementačnú časť práce, kde budú definované konkrétne API kontrakty, dátové modely, \glslink{dockercompose}{Docker Compose} konfigurácie a testovacie scenáre.